\chapter{Endoscopic Sinus Surgery}

\section*{What Is This Surgery?}
Endoscopic sinus surgery, most commonly referred to as functional endoscopic sinus surgery (FESS), is a minimally invasive surgical procedure performed through the nostrils to treat chronic inflammatory and infectious conditions of the paranasal sinuses. This procedure utilizes a rigid or flexible endoscope to provide magnified, illuminated visualization of the intricate anatomy of the nasal cavity and sinuses, allowing the surgeon to precisely identify and remove diseased tissue, polyps, and bone that obstruct the natural sinus drainage pathways. The primary goal of FESS is to restore normal mucociliary clearance and ventilation of the sinuses, thereby alleviating symptoms of chronic rhinosinusitis and preventing recurrent infections. It represents a significant advancement over older, more invasive external approaches, emphasizing the preservation of healthy mucosa and anatomical structures to promote physiological healing.

\section*{Alternative Names}
\begin{itemize}
  \item Functional Endoscopic Sinus Surgery (FESS)
  \item Endoscopic Sinusotomy
  \item Endoscopic Polypectomy
  \item Endoscopic Ethmoidectomy
  \item Endoscopic Maxillary Antrostomy
  \item Endoscopic Sphenoidotomy
  \item Endoscopic Frontal Sinusotomy
\end{itemize}

\section*{Relevant Surgical Anatomy}
A thorough understanding of the complex anatomy of the paranasal sinuses and their relationship to vital structures is paramount for safe and effective endoscopic sinus surgery. The paranasal sinuses consist of four paired air-filled cavities: the maxillary, ethmoid, frontal, and sphenoid sinuses. The maxillary sinuses are the largest, located within the maxilla, draining into the middle meatus via the maxillary ostium. The ethmoid sinuses are a labyrinth of 3-18 small air cells divided into anterior and posterior groups, separated by the basal lamella. They drain into the middle and superior meatus. The frontal sinuses, located in the frontal bone, drain into the frontal recess, which is part of the anterior ethmoid complex. The sphenoid sinuses, situated within the sphenoid bone at the skull base, drain into the sphenoethmoidal recess.

The ostiomeatal complex (OMC) is a critical anatomical region in the middle meatus where the maxillary, frontal, and anterior ethmoid sinuses drain. Obstruction of the OMC, often by the uncinate process, ethmoid bulla, or nasal polyps, is a common cause of chronic rhinosinusitis. Key anatomical landmarks and structures to consider during FESS include the lamina papyracea (the thin medial wall of the orbit, separating the ethmoid sinuses from the orbit), the skull base (roof of the ethmoid and sphenoid sinuses, separating them from the anterior cranial fossa), the optic nerve (in close proximity to the sphenoid sinus and posterior ethmoid cells), and the internal carotid artery (lateral to the sphenoid sinus). The anterior and posterior ethmoidal arteries, branches of the ophthalmic artery, traverse the ethmoid roof and are significant sources of bleeding if injured. The infraorbital nerve, a branch of the trigeminal nerve, exits the maxilla superior to the maxillary sinus and can be affected by extensive maxillary antrostomy. Lymphatic drainage of the sinuses primarily flows to the deep cervical lymph nodes.

\section{Surgical Principle \& Rationale}
The fundamental principle of FESS is the restoration of normal sinus physiology by addressing the underlying anatomical and pathological factors that impede mucociliary clearance and ventilation. The paranasal sinuses drain into the nasal cavity through small openings called ostia, which are often located within a critical anatomical region known as the ostiomeatal complex. Inflammation, mucosal edema, polyps, and anatomical variations (e.g., concha bullosa, deviated septum) can lead to obstruction of these ostia, trapping mucus and creating an environment conducive to bacterial or fungal growth, resulting in chronic rhinosinusitis.

FESS aims to enlarge these natural drainage pathways and remove obstructing pathology while meticulously preserving the healthy sinus mucosa and its ciliary function. The endoscopic approach provides magnified, illuminated visualization, allowing for precise dissection and targeted removal of diseased tissue. Key technical goals include:
\begin{itemize}
  \item Opening and enlarging the natural ostia of the maxillary, ethmoid, frontal, and sphenoid sinuses.
  \item Removing inflammatory polyps, hypertrophic mucosa, and other obstructing soft tissues.
  \item Correcting anatomical anomalies (e.g., septal spurs, concha bullosa) that contribute to obstruction.
  \item Minimizing trauma to healthy surrounding tissues and preserving the mucociliary blanket.
  \item Facilitating the delivery and efficacy of topical medications postoperatively by improving access to the sinus cavities.
\end{itemize}
By re-establishing proper ventilation and drainage, FESS helps to break the cycle of inflammation and infection, allowing the sinus mucosa to heal and function normally, thereby improving patient symptoms and quality of life.

\section*{History \& Surgical Evolution}
The concept of endoscopic visualization for nasal and sinus pathology dates back to the early 20th century with the work of Alfred Hirschmann, who developed the first endoscope for nasal examination in 1903. However, these early instruments were primarily diagnostic. The true revolution in sinus surgery began in the 1970s and 1980s, largely spearheaded by Professor Walter Messerklinger of Graz, Austria. Messerklinger conducted extensive anatomical and physiological studies of the paranasal sinuses, elucidating the importance of the ostiomeatal complex and the patterns of mucociliary flow. His work fundamentally shifted the understanding of sinusitis from a pan-sinus disease to one often originating from obstruction of key drainage pathways.

Building upon Messerklinger's insights, Professor Heinz Stammberger, also from Graz, developed and popularized the systematic endoscopic surgical techniques that focused on restoring function rather than radical exenteration, which characterized older external approaches like the Caldwell-Luc procedure. Their pioneering efforts, including detailed anatomical dissections and clinical application, led to the widespread adoption of FESS globally. The introduction of specialized endoscopic instruments, high-definition cameras, and subsequently, image-guided navigation systems in the late 20th and early 21st centuries, further enhanced the precision and safety of FESS, solidifying its position as the gold standard surgical treatment for chronic rhinosinusitis.

\section*{Clinical Indications}
Functional Endoscopic Sinus Surgery is indicated for a range of sinonasal pathologies, primarily when medical management has proven insufficient.

\begin{itemize}
  \item Chronic rhinosinusitis (CRS) with or without nasal polyps that is refractory to maximal medical therapy (e.g., antibiotics, topical corticosteroids, saline irrigations) for at least 12 weeks.
  \item Recurrent acute rhinosinusitis (RARS) with documented anatomical obstruction or persistent infection despite appropriate medical management.
  \item Sinus mucoceles or pyoceles, which are mucus-filled or pus-filled cysts that can expand and erode bone, requiring drainage and marsupialization.
  \item Fungal sinusitis, including allergic fungal rhinosinusitis (AFRS) or invasive fungal sinusitis, necessitating debridement and ventilation.
  \item Benign or malignant sinonasal tumors requiring endoscopic resection, biopsy, or debulking.
  \item Cerebrospinal fluid (CSF) leaks originating from the skull base within the sinonasal cavity, requiring endoscopic repair.
  \item Orbital decompression for thyroid eye disease or other causes of proptosis, to relieve pressure on the optic nerve.
  \item Endoscopic dacryocystorhinostomy (DCR) for lacrimal duct obstruction, to create a new drainage pathway for tears.
  \item Access for pituitary surgery or other anterior skull base procedures, providing a minimally invasive corridor.
\end{itemize}

\section*{Clinical Positioning}
Functional Endoscopic Sinus Surgery is generally considered a secondary treatment option for chronic rhinosinusitis, meaning it is typically reserved for patients who have failed an adequate trial of maximal medical therapy. This initial medical management often includes prolonged courses of antibiotics, oral or topical corticosteroids, nasal saline irrigations, and sometimes antifungal agents or leukotriene modifiers, usually for a minimum of 8-12 weeks. The role of FESS in these cases is to improve sinus ventilation and drainage, thereby enhancing the efficacy of ongoing topical medical therapies.

However, FESS can be considered a primary procedure in specific circumstances where immediate surgical intervention is necessary to address the pathology directly. Examples include the presence of a large, symptomatic mucocele, certain sinonasal tumors, or a cerebrospinal fluid leak, where delaying surgery could lead to significant morbidity or mortality. For most cases of chronic rhinosinusitis, the surgical intervention serves as an adjunct to ongoing medical management, aiming to improve the efficacy of topical medications by opening the sinus cavities and facilitating their delivery. It is rarely a standalone cure but rather a crucial component of a comprehensive, multidisciplinary treatment strategy.

\section*{Surgical Technique \& Procedure}
FESS is typically performed under general anesthesia, often in an outpatient setting, though some complex cases may require an overnight hospital stay. The patient is positioned supine with the head slightly elevated and extended. The procedure begins with the application of topical vasoconstrictors (e.g., oxymetazoline, epinephrine-soaked pledgets) to the nasal mucosa to reduce bleeding and improve visualization. Local anesthetic with epinephrine is then injected into key areas (e.g., greater palatine foramen, sphenopalatine ganglion region, uncinate process) to provide additional hemostasis and pain control.

The surgeon inserts a rigid endoscope (typically 0-degree or 30-degree) into the nasal cavity, providing a magnified, illuminated view of the anatomy. The specific steps vary depending on the extent of disease and the sinuses involved, but common operative steps include:

\begin{itemize}
  \item \textbf{Uncinectomy:} The uncinate process, a thin, sickle-shaped bone forming part of the lateral nasal wall, is carefully removed. This exposes the natural opening of the maxillary sinus (maxillary ostium) and the ethmoid bulla, which are key components of the ostiomeatal complex.
  \item \textbf{Maxillary Antrostomy:} The natural maxillary ostium is identified and enlarged, typically posteriorly, to improve drainage and ventilation of the maxillary sinus. This is often performed using a curved suction dissector or microdebrider.
  \item \textbf{Ethmoidectomy:} The ethmoid air cells, a labyrinth of small sinuses, are systematically opened and diseased tissue is removed. This can involve anterior ethmoidectomy (opening cells anterior to the basal lamella) and posterior ethmoidectomy (opening cells posterior to the basal lamella), with careful attention to the skull base and lamina papyracea.
  \item \textbf{Sphenoidotomy:} The sphenoid sinus ostium, located in the sphenoethmoidal recess in the posterior nasal cavity, is identified and enlarged to drain the sphenoid sinus. This requires careful identification of the superior turbinate and the sphenoid face.
  \item \textbf{Frontal Sinusotomy:} Access to the frontal sinus is achieved by opening the frontal recess, often requiring removal of agger nasi cells, frontal bulla cells, and other obstructing structures. For more extensive disease, a Draf procedure (e.g., Draf IIa, Draf III) may be performed to create a wider, more permanent opening.
  \item \textbf{Polypectomy:} Nasal polyps, if present, are carefully removed using microdebriders or forceps, ensuring complete removal from the sinus cavities.
\end{itemize}

Image-guided navigation systems, which use preoperative CT scans to create a real-time "GPS" for the surgeon, are frequently employed, especially in revision cases, extensive disease, or when operating near critical structures like the orbit or skull base. After the diseased tissue is removed and the drainage pathways are opened, the surgical site is meticulously inspected for hemostasis. Absorbable or non-absorbable packing may be placed in the nasal cavity, though minimal packing is preferred to reduce postoperative discomfort and facilitate early postoperative care. The nasal mucosa is not typically closed with sutures.

\section*{Preoperative Workup \& Preparation}
Thorough preoperative preparation is crucial for optimizing surgical outcomes and minimizing risks associated with FESS.

\begin{itemize}
  \item \textbf{Medical Optimization:} Patients should undergo a comprehensive medical evaluation to ensure they are fit for general anesthesia. This includes a detailed review of comorbidities (e.g., asthma, cardiac disease), current medications, and allergies. Any systemic conditions should be optimally managed prior to surgery.
  \item \textbf{Imaging:} A recent computed tomography (CT) scan of the paranasal sinuses, performed without intravenous contrast, is essential. This imaging provides a detailed anatomical roadmap for the surgeon, identifies the extent of disease, and highlights critical anatomical variants or areas of concern (e.g., dehiscence of the lamina papyracea, low-lying skull base, Onodi cells).
  \item \textbf{Maximal Medical Therapy:} For chronic rhinosinusitis, a documented trial of maximal medical therapy (e.g., 4-6 weeks of appropriate antibiotics, oral steroids, topical nasal steroids, saline irrigations) is usually required before surgery is considered, as per clinical guidelines.
  \item \textbf{Medication Adjustments:} Patients must discontinue blood-thinning medications (e.g., aspirin, NSAIDs, warfarin, novel oral anticoagulants) for a specified period before surgery, typically 5-10 days, as advised by the surgeon and primary care physician or cardiologist, to minimize bleeding risk.
  \item \textbf{NPO Status:} Patients must be NPO (nil per os - nothing by mouth) for 6-8 hours prior to surgery to prevent aspiration during anesthesia.
  \item \textbf{Prophylactic Antibiotics:} A single dose of intravenous prophylactic antibiotics (e.g., cefazolin) is typically administered preoperatively to reduce the risk of surgical site infection.
  \item \textbf{Informed Consent:} A detailed discussion of the procedure, its expected benefits, potential risks, and available alternatives is conducted, and informed consent is obtained from the patient.
\end{itemize}

\section*{Contraindications \& Precautions}
While FESS is a widely performed and generally safe procedure, certain conditions may contraindicate or necessitate extreme caution during surgery.

\textbf{Absolute Contraindications:}
\begin{itemize}
  \item Uncontrolled systemic disease (e.g., severe uncontrolled hypertension, unstable angina, recent myocardial infarction, severe coagulopathy) that poses an unacceptable anesthetic or surgical risk.
  \item Active, severe infection (e.g., acute orbital cellulitis with abscess, intracranial abscess) that requires immediate medical stabilization or alternative surgical approaches before elective sinus surgery.
  \item Patient refusal or inability to provide informed consent.
\end{itemize}

\textbf{Relative Contraindications \& Precautions:}
\begin{itemize}
  \item \textbf{Bleeding Disorders:} Uncontrolled coagulopathy or inability to safely discontinue anticoagulant medications significantly increases the risk of intraoperative and postoperative hemorrhage. Careful hematological evaluation and management are essential.
  \item \textbf{Severe Anatomic Distortion:} Extensive scarring from previous surgeries, severe polyposis obscuring landmarks, or significant anatomical variants (e.g., severely deviated septum, extensive concha bullosa) can make the surgery technically challenging and increase the risk of complications. Image-guided navigation is particularly important in these cases.
  \item \textbf{Immunocompromised State:} Patients who are severely immunocompromised (e.g., HIV/AIDS, organ transplant recipients on immunosuppressants) may have a higher risk of postoperative infection and impaired healing. Aggressive medical optimization is required.
  \item \textbf{Pediatric Patients:} While FESS can be performed in children, the indications are more stringent, and careful consideration of facial growth and development, as well as the smaller anatomical spaces, is necessary.
  \item \textbf{Patient Compliance:} Poor patient compliance with postoperative care instructions, including nasal irrigations and follow-up appointments, can negatively impact long-term outcomes.
\end{itemize}
The surgeon must exercise extreme caution due to the proximity of vital structures such as the orbit (eye), optic nerve, and anterior cranial fossa (brain). The thin lamina papyracea (medial orbital wall) and the skull base are particularly vulnerable to injury.

\section*{Complications \& Risks}
While FESS is generally safe, as with any surgical procedure, there are potential risks and complications. Serious complications are rare but can occur.

\begin{itemize}
  \item \textbf{General Surgical Complications:}
    \begin{itemize}
      \item \textbf{Bleeding (Epistaxis):} Common during and after surgery. Rarely severe enough to require blood transfusion, nasal packing, or re-operation.
      \item \textbf{Infection:} Postoperative sinusitis or wound infection, typically managed with antibiotics.
      \item \textbf{Anesthesia Complications:} Risks associated with general anesthesia, including adverse drug reactions, respiratory issues, or cardiovascular events.
    \end{itemize}
  \item \textbf{Procedure-Specific Risks:}
    \begin{itemize}
      \item \textbf{Adhesion Formation (Synechiae):} Scar tissue can form between mucosal surfaces within the nasal cavity or sinuses, potentially re-obstructing the sinus ostia. This may require in-office lysis or revision surgery.
      \item \textbf{Recurrence of Disease:} Chronic rhinosinusitis, especially with nasal polyps, can recur despite successful surgery, necessitating ongoing medical therapy or revision surgery.
      \item \textbf{Orbital Injury:}
        \begin{itemize}
          \item \textbf{Periorbital Ecchymosis/Edema:} Bruising and swelling around the eye, usually temporary and self-resolving.
          \item \textbf{Diplopia (Double Vision):} Due to injury to extraocular muscles or orbital fat, often transient, but can be permanent in rare cases.
          \item \textbf{Orbital Hematoma:} Bleeding into the orbit, which can cause proptosis (bulging eye) and, if severe, compress the optic nerve leading to vision loss. Requires urgent intervention.
          \item \textbf{Blindness:} Extremely rare, but possible due to direct optic nerve injury or severe orbital hematoma with optic nerve compression.
        \end{itemize}
      \item \textbf{Cerebrospinal Fluid (CSF) Leak:} Injury to the skull base can create an opening between the brain and the nasal cavity, leading to CSF leakage. This requires prompt repair to prevent meningitis and other intracranial complications.
      \item \textbf{Anosmia/Hyposmia:} Loss or reduction of the sense of smell, which can be temporary due to swelling or permanent due to direct injury to the olfactory neuroepithelium.
      \item \textbf{Numbness:} Numbness in the upper front teeth, cheek, or palate due to injury to the infraorbital nerve, usually temporary.
      \item \textbf{Septal Perforation:} A hole in the nasal septum, which can cause crusting, bleeding, or a whistling sound during breathing.
      \item \textbf{Dental Injury:} Damage to teeth during intubation or manipulation of instruments.
    \end{itemize}
\end{itemize}

\section*{Postoperative Recovery Timeline}
Immediately after FESS, patients are transferred to the Post-Anesthesia Care Unit (PACU) for close monitoring as they recover from general anesthesia. Pain is typically managed with oral analgesics, and some nasal congestion, mild bloody drainage, and facial pressure are common. Most patients are discharged home the same day, provided they are stable, have adequate pain control, and no immediate complications. Patients are advised to rest with their head elevated, avoid strenuous activities, and refrain from blowing their nose forcefully for the first 24-48 hours.

During the first 1-2 weeks, nasal saline irrigations, often starting within 24-48 hours, are crucial for clearing blood clots, crusting, and mucus, promoting healing, and preventing adhesions. Oral antibiotics and corticosteroids may be prescribed for a short course to reduce infection risk and inflammation. Patients may experience continued nasal congestion, mild fatigue, and occasional bloody discharge. The surgeon will typically schedule a follow-up visit within this period for endoscopic debridement, where any remaining clots, crusts, or early adhesions are gently removed. Long-term recovery involves continued nasal hygiene, often with topical nasal steroids, and regular follow-up appointments to monitor healing and manage any recurrence of symptoms or polyps. Full mucosal healing and stabilization of symptoms can take several weeks to months, with patients gradually resuming normal activities over this period.

\section*{Surgical Findings \& Pathology}
During FESS, the surgeon meticulously documents the intraoperative findings, which typically include the extent of mucosal inflammation, the presence and size of nasal polyps, the patency of sinus ostia, and any anatomical variants contributing to obstruction (e.g., concha bullosa, septal deviation, paradoxical middle turbinate). The specific sinuses involved and the degree of disease within each are also noted. Any tissue removed during the procedure, such as polyps, hypertrophic mucosa, or bone fragments, is sent for histopathological examination. The pathology report is a critical component of the patient's diagnosis and management. It will confirm the nature of the disease, such as chronic inflammation, the presence of specific types of polyps (e.g., eosinophilic polyps, fungal polyps), fungal elements (e.g., mucin consistent with allergic fungal rhinosinusitis), or, rarely, a benign or malignant tumor. This report guides further medical therapy and helps predict the long-term prognosis.

\section*{Expected Postoperative Course}
A normal, successful postoperative course following FESS is characterized by a gradual and sustained improvement in the patient's preoperative symptoms. Initially, patients may experience nasal congestion, mild pain, and some bloody discharge, which progressively diminish over the first few weeks. Pain should be manageable with oral analgesics and resolve within days. As healing progresses, endoscopic examination reveals open, patent sinus ostia with a healthy, re-epithelialized mucosal lining that is free from significant inflammation, crusting, or synechiae (adhesions). The nasal passages should appear clear and well-ventilated. Clinically, the patient experiences a significant reduction in nasal congestion, decreased facial pain and pressure, less purulent discharge, and often an improvement in the sense of smell. The frequency and severity of acute exacerbations of sinusitis should decrease substantially, leading to an overall enhanced quality of life and reduced reliance on antibiotics.

\section{Postoperative Care \& Follow-up}
Postoperative care is critical for optimizing the long-term success of FESS and preventing recurrence.

\begin{itemize}
  \item \textbf{Postoperative Endoscopic Debridement:} Patients typically require several in-office endoscopic debridement procedures in the weeks to months following surgery. During these visits, the surgeon removes crusts, blood clots, and any early adhesions to ensure the sinus ostia remain open and to promote proper mucosal healing.
  \item \textbf{Continued Medical Therapy:} Long-term medical management is essential to maintain surgical benefits. This often includes daily nasal saline irrigations, topical nasal steroid sprays, and sometimes other medications like leukotriene inhibitors or biologics, especially for patients with nasal polyps or allergic fungal rhinosinusitis.
  \item \textbf{Activity Resumption:} Patients are advised to gradually resume normal activities, avoiding strenuous exercise, heavy lifting, and nose blowing for the first few weeks. Air travel is generally safe after a few days, but diving should be avoided for several weeks.
  \item \textbf{Imaging Surveillance:} In some cases, particularly for extensive disease, recurrent polyps, or suspected complications, follow-up CT scans may be performed several months after surgery to assess the anatomical results and identify any recurrence.
  \item \textbf{Warning Signs:} Patients are instructed to contact their surgeon immediately if they experience severe, uncontrolled bleeding, worsening pain, fever, vision changes (e.g., double vision, decreased acuity), severe headache, or clear watery nasal discharge (suggesting a CSF leak).
\end{itemize}
Regular follow-up appointments are crucial for monitoring the long-term health of the sinuses and adjusting medical therapy as needed to prevent recurrence and optimize outcomes.

\section*{Epidemiology}
Chronic rhinosinusitis (CRS), the primary indication for FESS, is a highly prevalent condition, affecting approximately 10-15\% of the adult population in developed countries. Its incidence varies globally but represents a significant public health burden due to its impact on quality of life, productivity, and healthcare costs. CRS can occur at any age but is most commonly diagnosed in adults between 30 and 60 years old, with no significant predilection for sex, although certain phenotypes like CRS with nasal polyps (CRSwNP) may show a slight male predominance. The frequency of specific indications for FESS largely mirrors the prevalence of CRS subtypes; CRSwNP accounts for a substantial proportion of surgical cases, often requiring revision surgery due to high recurrence rates. Other indications, such as mucoceles or fungal sinusitis, are less common. While FESS is a very safe procedure, minor morbidity such as postoperative bleeding or synechiae formation is relatively common, occurring in 5-15\% of cases. Serious complications, including orbital injury, CSF leak, or blindness, are exceedingly rare, with reported rates typically less than 1\% and often closer to 0.1-0.5\%. Mortality directly attributable to FESS is virtually non-existent.

\section*{Outcomes \& Prognosis}
The outcomes of FESS are generally favorable, with typical success rates for significant symptom improvement ranging from 80\% to 90\% in carefully selected patients. Functional outcomes, such as improved nasal breathing, reduced facial pain, and enhanced quality of life, are consistently reported across numerous studies. The prognosis is influenced by several factors, including the underlying disease phenotype (e.g., CRS without polyps generally has a better prognosis than CRSwNP), the presence of comorbidities (e.g., asthma, aspirin-exacerbated respiratory disease, allergic rhinitis), and patient adherence to postoperative medical therapy. For patients with nasal polyps, recurrence rates can be as high as 40-60\% within five years, often necessitating revision surgery and aggressive long-term medical management, including topical steroids and potentially biologics. However, even with recurrence, subsequent surgeries are often less extensive, and symptoms are typically less severe than the initial presentation. Overall, FESS significantly reduces the burden of chronic rhinosinusitis, allowing most patients to achieve better disease control and a substantial improvement in their daily lives. Landmark studies, such as the International Consensus Statement on Allergy and Rhinology: Rhinosinusitis (ICAR:RS), consistently support FESS as an effective treatment when medical therapy fails.

\section*{References}
\begin{enumerate}
  \item MedlinePlus. Functional endoscopic sinus surgery. U.S. National Library of Medicine; 2024. Available from: \url{https://medlineplus.gov/ency/article/007204.htm}
  \item Lal D, Scangas GA, Bleier BS. Functional Endoscopic Sinus Surgery. In: Flint PW, Haughey BH, Lund VJ, et al., editors. Cummings Otolaryngology: Head and Neck Surgery. 7th ed. Philadelphia, PA: Elsevier; 2021. p. 667-680.
  \item Orlandi RR, Kingdom TT, Hwang PH, et al. International Consensus Statement on Allergy and Rhinology: Rhinosinusitis 2021. Int Forum Allergy Rhinol. 2021 Mar;11(3):213-739.
  \item Rosenfeld RM, Piccirillo JF, Chandrasekhar SS, et al. Clinical Practice Guideline: Adult Sinusitis. Otolaryngol Head Neck Surg. 2015 Apr;152(1 Suppl):S1-S39.
\end{enumerate}

\clearpage